\documentclass[12pt]{article}

\usepackage{amsmath}
\usepackage{enumitem}
\usepackage{geometry}
\usepackage{titlesec}

% Page settings
\geometry{
    a4paper,
    margin=0.75in
}

% Reduce space before and after sections
\titleformat{\section}
  {\large\bfseries}
  {\thesection}
  {0.5em}
  {}

\titlespacing*{\section}
  {0pt}
  {10pt}
  {4pt}

% Reduce paragraph spacing
\setlength{\parindent}{0pt}
\setlength{\parskip}{4pt}

\begin{document}

\begin{center}
    {\LARGE \textbf{Prompting Techniques}}
\end{center}

\vspace{8pt}

\section{Zero-Shot Prompting}

Asking the model to do a task with no examples given.

\textbf{Example:}

\begin{quote}
``Calculate the area of a rectangle with length 12 cm and width 5 cm.''
\end{quote}


\section{Few-Shot Prompting}

Giving a few examples before the actual task so the model can learn the pattern.

\textbf{Example:}

\begin{quote}
English: Good morning $\rightarrow$ French: Bonjour

English: Thank you $\rightarrow$ French: Merci

English: See you soon $\rightarrow$ French: ?
\end{quote}


\section{Chain-of-Thought (CoT) Prompting}

Asking the model to reason step-by-step before giving the final answer.

\textbf{Example:}

\begin{quote}
``A student has 25 chocolates and gives 8 to a friend. Then the student receives 10 more chocolates. How many chocolates does the student have? Think step by step.''
\end{quote}


\section{Zero-Shot Chain-of-Thought}

Just adding ``think step by step'' without giving examples.

\textbf{Example:}

\begin{quote}
``A shirt costs ₹800 and has a 25\% discount. What is the final price? Let's think step by step.''
\end{quote}


\section{Role-Based (Persona) Prompting}

Assigning a role or identity to the model to shape its tone/expertise.

\textbf{Example:}

\begin{quote}
``You are an experienced Python teacher. Explain what a dictionary is in Python to a beginner.''
\end{quote}


\section{Instruction-Based Prompting}

Giving clear, direct instructions about the desired output format or behavior.

\textbf{Example:}

\begin{quote}
``Explain the advantages of artificial intelligence in exactly 5 bullet points.''
\end{quote}


\section{Contextual Prompting}

Providing background/context before asking the question.

\textbf{Example:}

\begin{quote}
``I'm a beginner in programming. Explain what a Python function is using simple language and an easy example.''
\end{quote}


\section{Self-Consistency Prompting}

Asking the model to generate multiple reasoning paths and pick the most consistent answer (usually done programmatically, not just in one prompt).

\textbf{Example:}

\begin{quote}
``Solve this math problem using 5 different approaches and choose the answer that appears most frequently.''
\end{quote}


\section{Prompt Chaining}

Breaking a complex task into multiple prompts, where the output of one feeds into the next.

\textbf{Example:}

\begin{quote}
\textbf{Prompt 1:} ``Extract the important points from this article.''

\textbf{Prompt 2:} ``Using those points, create a short summary.''

\textbf{Prompt 3:} ``Convert the summary into 5 bullet points.''
\end{quote}


\section{Tree-of-Thought Prompting}

The model explores multiple possible reasoning branches before settling on an answer (more advanced/structured than CoT).

\textbf{Example:}

\begin{quote}
``Consider three different ways to solve this problem. Compare each method and select the most effective one.''
\end{quote}


\section{Retrieval-Augmented Prompting (RAG-style)}

Providing external documents/data as context, then asking questions based on that content.

\textbf{Example:}

\begin{quote}
``Here is a document about a company's leave policy: [text]. Based on this document, how many paid leaves are available to employees?''
\end{quote}


\section{Negative Prompting}

Telling the model what NOT to do or include.

\textbf{Example:}

\begin{quote}
``Write a short explanation of machine learning, but do not use technical terms or complicated language.''
\end{quote}

\end{document}
